% ============================================================================
\clearpage
\section{Journey Map del equipo}\label{sec:journey}

\subsection{Objetivo mapeado y persona seleccionada}

Tras discutir los seis escenarios, el equipo seleccionó el objetivo de Laura Méndez:
\textbf{validar una cifra de cartera de crédito y responder con la fuente citada antes de una
reunión}. La elección se apoyó en tres criterios:

\begin{uxlist}
\lead{Persona primaria}{la actividad anterior definió a Laura como el blanco de diseño
  específico del producto.}
\lead{Tarea compartida}{los demás perfiles también localizan o validan datos, aunque esa no
  sea siempre su actividad principal.}
\lead{Variedad de contactos}{el recorrido vincula buscador, catálogo, filtros, asistente,
  consultas guardadas y avisos de versión.}
\end{uxlist}

El mapa comienza con la solicitud y termina cuando Laura revisa, dos semanas después, un aviso
de cambio. Así cubre los momentos antes, durante y después del uso sin ampliar el alcance a
otras tareas.

\subsection{Listado de actividades}

El equipo descompuso el objetivo en las actividades que alimentan el carril de acciones.

\begin{uxtabla}{C{1.1cm}L{3.2cm}Z{1.0}}{Listado de actividades que Laura ejecuta para cumplir el objetivo}
\uxheadrow \thd{\#} & \thd{Actividad} & \thd{Qué implica} \\ \midrule
1 & Recibir e interpretar la solicitud & Identificar el concepto, el corte y el plazo \\
2 & Decidir dónde buscar & Elegir entre el portal, un archivo anterior o preguntar a un colega \\
3 & Acceder al portal & Iniciar sesión y llegar al espacio de trabajo de su rol \\
4 & Formular la búsqueda & Escribir el concepto en lenguaje de negocio \\
5 & Evaluar los resultados & Distinguir entre variables de nombre parecido \\
6 & Descartar versiones retiradas & Reconocer cuál es la vigente \\
7 & Leer la definición oficial & Confirmar que la variable es la que le piden \\
8 & Identificar al propietario & Saber a quién acudir si necesita respaldo \\
9 & Verificar la vigencia & Confirmar la fecha de última actualización \\
10 & Acotar al periodo solicitado & Aplicar el filtro sin perder el contexto \\
11 & Resolver una duda de alcance & Preguntar al asistente y verificar la fuente de la respuesta \\
12 & Obtener el dato con su procedencia & Copiar valor, fuente y fecha en un solo paso \\
13 & Responder la solicitud & Redactar y enviar con la referencia incluida \\
14 & Guardar la consulta & Nombrarla para poder reutilizarla \\
15 & Atender el cambio posterior & Revisar el aviso de nueva versión y su impacto \\
\end{uxtabla}

\subsection{El taller de construcción}\label{sec:taller}

Los tres integrantes construyeron el mapa en una sesión remota de dos horas. La plantilla
contenía los carriles vacíos; durante el taller se acordaron primero las etapas y luego las
acciones, contactos, pensamientos, emociones, fricciones y oportunidades. Las dudas se
marcaron como hipótesis para validación posterior.

\begin{uxtabla}{L{3.0cm}Z{1.0}Z{1.0}}{Síntesis del taller de journey mapping}
\uxheadrow \thd{Elemento} & \thd{Decisión} & \thd{Propósito} \\ \midrule
Objetivo & Mapear la validación urgente de Laura & Mantener un alcance común y verificable \\
Participantes & Los tres integrantes con sus escenarios y personas & Integrar perspectivas distintas del producto \\
Materiales & Videoconferencia y pizarra digital con carriles preparados & Concentrar la sesión en el contenido \\
Secuencia & Etapas, experiencia actual, fricciones y oportunidades & Evitar proponer soluciones antes de comprender el recorrido \\
Cierre & Registro de supuestos y del instrumento que los validará & Vincular el mapa con la investigación posterior \\
\end{uxtabla}

El taller incorporó la etapa de reutilización, que no estaba prevista al inicio. Al recorrer
la experiencia completa, el equipo observó que una definición puede cambiar después de la
entrega y afectar el trabajo ya realizado.

\subsection{Estructura del mapa}

Las quince actividades se agrupan en \textbf{siete etapas} y \textbf{siete carriles}. Los
elementos centrales de la consigna se complementan con puntos de contacto y una escala
emocional para facilitar la lectura.

\begin{uxtabla}{L{3.3cm}Z{1.0}Z{1.0}}{Carriles del journey map y su correspondencia con lo solicitado}
\uxheadrow \thd{Carril} & \thd{Qué documenta} & \thd{Elemento solicitado} \\ \midrule
Etapas & Los momentos en que se divide el recorrido, en orden cronológico & Cómo el usuario aborda cada etapa \\
Acciones & Lo que el usuario hace en cada etapa & Cómo el usuario aborda cada etapa \\
Pensamientos & Lo que se pregunta o se dice a sí mismo & Sus pensamientos \\
Pain points & Lo que complica la tarea y puede impedir la adopción del producto & Los puntos que pueden complicar la adopción \\
Oportunidades & Las decisiones de diseño que mitigan cada punto de fricción & Las oportunidades que pueden mitigarlos \\
Puntos de contacto & Con qué parte del producto o de su entorno interactúa & Complemento del ejemplo de referencia \\
Emociones & Cómo se siente, en una escala declarada de $-2$ a $+2$ & Complemento del ejemplo de referencia \\
\end{uxtabla}

Los pensamientos y emociones retoman los cuadrantes \emph{Thinks} y \emph{Feels} del mapa de
empatía de Laura. De esta manera, el journey map conserva trazabilidad con la actividad
anterior.

\begin{uxnota}[Herramienta de construcción]
Los cuatro mapas de este documento conservan la estructura del ejemplo de referencia de la
actividad, elaborado con la herramienta propuesta: la misma retícula de etapas por carriles,
con la curva emocional entre los pensamientos y los puntos de contacto. La composición final
se produjo con una plantilla propia del equipo que imprime esa retícula en PDF vectorial, por
dos razones. La primera es de legibilidad: a esta densidad de texto, un mapa exportado como
imagen se pixela en cuanto el lector amplía una tarjeta, mientras que el trazado vectorial
conserva el detalle a cualquier escala. La segunda es de consistencia: los cuatro mapas
comparten así la tipografía, la paleta y la escala emocional del resto del documento, y los
cuatro se generan desde una única fuente de contenido, de modo que corregir una celda no
puede dejar una versión desactualizada en otro mapa.
\end{uxnota}

\subsection{El mapa}

La vista panorámica se presenta en orientación horizontal para conservar la lectura conjunta
de las siete etapas.

\figurapanoramica{figuras/journey_equipo_laura.pdf}
  {Journey map de equipo. Persona: Laura Méndez. Objetivo: validar una cifra de cartera de
   crédito y responder con la fuente citada antes de una reunión. Elaborado en taller por los
   tres integrantes.}{fig:journey}

\subsection{Lectura del recorrido}\label{sec:recorrido}

Antes de entrar al portal, Laura depende de su memoria y de sus contactos para decidir dónde
buscar. El acceso debe llevarla directamente a una búsqueda en lenguaje de negocio; después,
las tarjetas necesitan distinguir variables similares mediante fuente, propietario y vigencia.

La validación es el \textbf{momento de la verdad}. La ficha debe reunir la definición oficial,
las notas del propietario y el historial de versiones. Si resuelve esas preguntas, Laura
mantiene el control; si necesita escribir a un especialista, la tarea vuelve a depender de un
tercero.

El recorrido termina con la entrega de una cifra que conserva su procedencia y con una
consulta reutilizable. El aviso posterior de cambio extiende el valor del portal más allá de
la tarea inmediata, porque permite revisar el impacto sobre resultados ya compartidos.

\begin{uxmomento}
\textbf{Momento de la verdad.} La relación con el producto se define en la validación: el
portal debe sustituir la dependencia del conocimiento informal, no convertirse en un paso
adicional antes de consultar a otra persona.
\end{uxmomento}

\subsection{Curva emocional}

\figuraux{figuras/curva_emocional.pdf}
  {Curva emocional del recorrido de Laura Méndez, en una escala declarada de $-2$ a $+2$. La
   línea continua corresponde al recorrido con el portal; la punteada, al recorrido actual sin
   él. Elaboración propia.}{fig:curva}

La escala representa bloqueo en $-2$, neutralidad en $0$ y confianza en $+2$. La diferencia
principal aparece en la validación: hoy la dependencia de un tercero detiene la tarea; con el
portal, la misma etapa alcanza el punto más alto al ofrecer contexto verificable.

\subsection{Brechas de conocimiento}

Las brechas siguientes indican qué necesita saber Laura para avanzar y dónde debe encontrarlo.

\begin{uxtabla}{C{1.6cm}Z{1.0}Z{1.0}}{Brechas de conocimiento por etapa}
\uxheadrow \thd{Etapa} & \thd{Qué necesita saber y no sabe} & \thd{Dónde debe resolverse} \\ \midrule
1. Disparo & Si el portal contiene el dato que le piden & Alcance del catálogo visible desde fuera del portal \\
3. Búsqueda & Cuál de dos variables similares corresponde a su solicitud & Diferenciación en la tarjeta de resultado \\
4. Validación & Si la definición cambió desde la última vez que la usó & Historial de versiones en la ficha \\
5. Profundización & De dónde obtuvo el asistente cada cifra & Visibilidad de la herramienta consultada y cita por cifra \\
6. Entrega & Qué debe acompañar al número para que sea defendible & Copiado con procedencia por omisión \\
7. Reutilización & Si lo que entregó antes del cambio sigue siendo válido & Aviso con alcance retroactivo declarado \\
\end{uxtabla}

\subsection{Recorridos entrelazados}\label{sec:entrelazados}

El recorrido de Laura depende de tareas realizadas por otros perfiles y también alimenta sus
decisiones.

\figuraux{figuras/journeys_entrelazados.pdf}
  {Puntos de encuentro entre el recorrido de Laura Méndez y los de las otras personas.
   Elaboración propia.}{fig:entrelazados}

\begin{uxlist}
\lead{Etapa 2 y escenario 5}{el acceso correcto depende del alta y del permiso mínimo que
  administra Mariana Ovalle.}
\lead{Etapas 4 y 7 con escenario 4}{la definición y el aviso de cambio proceden del trabajo de
  gobierno que realiza Roberto Valdez.}
\lead{Etapa 6 y escenario 3}{la cifra validada puede alimentar el indicador que Arturo
  Castañeda consulta antes de la Junta.}
\lead{Etapa 3 y escenario 2}{el descubrimiento puntual de fuentes relacionadas corresponde al
  modo de trabajo habitual de Diego Hernández.}
\end{uxlist}

Estas relaciones muestran que algunas fricciones no se resuelven solo en la interfaz de
Laura. La calidad de la validación y de los avisos también depende de los procesos de gobierno
y administración que sostienen el portal.

